\documentclass[../linear-algebra.tex]{subfiles}

\setcounter{MaxMatrixCols}{20}

\begin{document}
\section{System of Linear Equations}
\subsection{Linear Equations}
A linear equation in $n$ variables is
\[a_1x_1 + a_2x_2 + \cdots + a_nx_n = b\]
The coefficients $a_1,\ldots,a_n$ cannot all be 0, 
and variables $x_1,\ldots,x_n$ are raised only to the first power.

\noindent A linear system consisting of $m$ linear equations in $n$ variables
$x_1, x_2, \cdots, x_n$, is represented as
\begin{alignat*}{4}
    a_{11}x_1 &+ a_{12}x_2 &+ \cdots &+ a_{1n}x_n &&= b_1 \\
    a_{21}x_1 &+ a_{22}x_2 &+ \cdots &+ a_{2n}x_n &&= b_2 \\
    \vdots    & \quad &         & \quad \vdots && \quad \vdots \\
    a_{m1}x_1 &+ a_{m2}x_2 &+ \cdots &+ a_{mn}x_n &&= b_m
\end{alignat*}

\noindent The solution of a linear system in $n$ unknowns is a sequence of $n$ numbers
$s_1, s_2,\ldots,s_n$.

\medskip

\noindent Any system of linear equations behave in one of these ways:
\begin{enumerate}
    \item The system has no solution
    \item The system has a unique solution
    \item The system has infintely many solutions
\end{enumerate}

\noindent If the linear system has at least one solution, it is \textbf{consistent}. 

\subsection{Homogenous Linear System}
\noindent If $b_1=b_2=\ldots=b_n=0$, it is a homogenous linear system.
Every homogenous system has at least one solution, known as the trivial solution,
\[x_1 = x_2 = \ldots = x_n = 0\]
Hence, every homogenous system is consistent, and 
\begin{enumerate}
    \item has only the trivial solution, or
    \item has infintely many solutions in addition to the trivial solution
\end{enumerate}

\subsection{Representation as Matrix}
\noindent Consider a linear system of two equations:
\begin{align*}
    5x + y &= 3 \\
    2x - y &= 4
\end{align*}

\noindent Coefficient matrix:
$\begin{bmatrix} 
    5 & 1 \\ 
    2 & -1 
\end{bmatrix}$

\medskip

\noindent Augmented matrix:
$\begin{bmatrix} 
    5 & 1 & 3 \\ 
    2 & -1 & 4 
\end{bmatrix}$

\subsection{The two fundamental questions about any linear system}
\begin{enumerate}
    \item Is the system consistent?
    \item If a solution exists, is it unique?
\end{enumerate}
The following concepts are tools to answer these questions.

\subsection{Row Echelon and Reduced Row Echelon Form}
Row equivalent matrices can be transformed to each other 
by one of these elementary row operations:
\begin{enumerate}
    \item \textbf{Swap}: Swapping two rows
    \item \textbf{Scale}: Multiply a row by a non-zero constant
    \item \textbf{Pivot}: Add a multiple of one row to another row
\end{enumerate}

\subsubsection{Row Echelon Form}
Properties:
\begin{enumerate}
    \item All rows having only zero entries are at the bottom.
    \item The leftmost non-zero entry (i.e. leading entry) of every non-zero row (i.e. the pivot)
    is to the right of the leading entry of every row above.
\end{enumerate}
A matrix can have multiple REF forms, depending on the steps taken.
Matrix in row echelon form (and not in reduced row echelon form):
\[\begin{bmatrix}
    1 & a_0 & a_1 & a_2 & a_3 \\
    0 & 0   & 2   & a_4 & a_5 \\
    0 & 0   & 0   & 1   & a_6 \\
    0 & 0   & 0   & 0   & 0
\end{bmatrix}\]

\subsubsection{Reduced Row Echelon Form}
Properties:
\begin{enumerate}
    \item The matrix is in row echelon form.
    \item The leading entry (or pivot) of every non-zero row is 1.
    \item Every column containing a leading 1 has zeroes in its other entries.
\end{enumerate}
Every matrix has a unique RREF matrix, regardless of the steps taken.

\smallskip

\noindent Example matrix in reduced row echelon form:
\[\begin{bmatrix}
    1 & 0 & 0 & a_1 & a_2 \\
    0 & 1 & 0 & a_3 & a_4 \\
    0 & 0 & 1 & a_5 & a_6 \\
    0 & 0 & 0 & 0   & 0
\end{bmatrix}\]

\subsubsection{Gaussian and Gauss-Jordan elimination}
Consider this linear of system equations:
\begin{alignat*}{4}
     2x &+  y &&-  z &&= &  8 \\
    -3x &-  y &&+ 2z &&= &-11 \\
    -2x &+  y &&+ 2z &&= &-3
\end{alignat*}
% The augmented matrix is:
% \[\begin{bmatrix}
%      2 &  1 & -1 &   8 \\
%     -3 & -1 &  2 & -11 \\
%     -2 &  1 &  2 &  -3
% \end{bmatrix}\]
% \[
% \begin{array}{r|c}
%     \begin{aligned}
%         L_2 + \frac{3}{2}L_1 &\rightarrow L_2 \\
%         L_3 + L_1 &\rightarrow L_3
%     \end{aligned} 
%     & 
%     \centering
%     \left[
%     \begin{array}{ccc|c}
%         2 & 1 & -1 & 8 \\
%         0 & \frac{1}{2} & \frac{1}{2} & 1 \\
%         0 & 2 & 1 & 5
%     \end{array}
%     \right]
% \end{array}
% \]
% \[
% \begin{array}{r|c}
%     L_3 + -4L_2 \rightarrow L_3 
%     & 
%     \left[
%     \begin{array}{ccc|c}
%         2 & 1 & -1 & 8 \\
%         0 & \frac{1}{2} & \frac{1}{2} & 1 \\
%         0 & 0 & -1 & 1
%     \end{array}
%     \right]
% \end{array}
% \]

\begin{table}[H] % [H] forces the table to stay HERE
\centering
\renewcommand{\arraystretch}{1.2} 

\begin{tabular}{|>{\centering\arraybackslash}m{4.5cm}|>{\centering\arraybackslash}m{6.5cm}|}
    \hline
    \textbf{Row operations} & \textbf{Augmented matrix} \\ \hline
    
    % --- Row 1 ---
    % \rule{0pt}{2.5em} creates invisible top padding (smaller than before)
    \rule{0pt}{2.5em} 
    & 
    $\left[
    \begin{array}{ccc|c}
         2 &  1 & -1 &   8 \\
        -3 & -1 &  2 & -11 \\
        -2 &  1 &  2 &  -3
    \end{array}
    \right]$ 
    \\[2em] \hline % [2em] creates invisible bottom padding
    
    % --- Row 2 ---
    \rule{0pt}{2.5em} 
    $\begin{aligned}
        L_2 + \tfrac{3}{2}L_1 &\rightarrow L_2 \\
        L_3 + L_1 &\rightarrow L_3
    \end{aligned}$ 
    & 
    $\left[
    \begin{array}{ccc|c}
        2 & 1 & -1 & 8 \\
        0 & \tfrac{1}{2} & \tfrac{1}{2} & 1 \\
        0 & 2 & 1 & 5
    \end{array}
    \right]$ 
    \\[2em] \hline
    
    % --- Row 3 ---
    \rule{0pt}{2.5em} 
    $L_3 + -4L_2 \rightarrow L_3$ 
    & 
    $\left[
    \begin{array}{ccc|c}
        2 & 1 & -1 & 8 \\
        0 & \frac{1}{2} & \frac{1}{2} & 1 \\
        0 & 0 & -1 & 1
    \end{array}
    \right]$ 
    \\[2em] \hline
\end{tabular}
\end{table}
\noindent We obtain the the row echelon form using Gaussian elimination.
Note that top-down approach uses the rows above to introduce zeroes below the pivots.

\noindent This is also the forward phase of Gauss-Jordan elimination.

\begin{table}[H]
\centering
\renewcommand{\arraystretch}{1.2}

\begin{tabular}{|>{\centering\arraybackslash}m{4.5cm}|>{\centering\arraybackslash}m{6.5cm}|}
    \hline
    
    \textbf{Row operations} & \textbf{Augmented matrix} \\ \hline
    % --- Step 1 ---
    \rule{0pt}{2.5em} 
    $\begin{aligned}
        L_1 - L_3 &\rightarrow L_1 \\
        L_2 + \frac{1}{2}L_3 &\rightarrow L_2
    \end{aligned}$ 
    & 
    $\left[
    \begin{array}{ccc|c}
         2 &  1 &  0 &  7 \\
         0 & \frac{1}{2} &  0 & \frac{3}{2} \\
         0 &  0 & -1 &  1
    \end{array}
    \right]$ 
    \\[2em] \hline
    
    % --- Step 2 ---
    \rule{0pt}{2.5em} 
    $\begin{aligned}
        2L_2 &\rightarrow L_2 \\
        -L_3 &\rightarrow L_3
    \end{aligned}$ 
    & 
    $\left[
    \begin{array}{ccc|c}
        2 & 1 & 0 & 7 \\
        0 & 1 & 0 & 3 \\
        0 & 0 & 1 & -1
    \end{array}
    \right]$ 
    \\[2em] \hline
    
    % --- Step 3 ---
    \rule{0pt}{2.5em} 
    $\begin{aligned}
        L_1 - L_2 &\rightarrow L_1 \\
        \frac{1}{2}L_1 &\rightarrow L_1
    \end{aligned}$ 
    & 
    $\left[
    \begin{array}{ccc|c}
        1 & 0 & 0 & 2 \\
        0 & 1 & 0 & 3 \\
        0 & 0 & 1 & -1
    \end{array}
    \right]$ 
    \\[2em] \hline
\end{tabular}
\end{table}
\noindent We obtain the RREF using Gauss-Jordan elimination.
Note that bottom-up approach uses the rows below to introduce zeroes above pivots,
and to transform the pivots to 1.

\noindent This is the backward phase of Gauss-Jordan elimination.

\medskip

\noindent Note that of the EROs, only Scale and Pivot were used.

\subsection{Answering the fundamental questions}
Using row reduction to solve a linear system:
\begin{enumerate}
    \item Write the augmented matrix of the linear system.
    \item Use row reduction algorithms to obtain row echelon form 
    of the augmented matrix. Here we can decide if the system in consistent.
    If it is, go to the next step.
    \item Continue row reduction algorithm to obtain reduced row echelon form.
    \item Write the system of equations corresponding to this matrix.
    \item Rewrite each nonzero equation so that the leading variable
    is expressed in terms of any free variables.    
\end{enumerate}

\subsubsection{Example (Not consistent system)}
Determine if this linear system is consistent:
\begin{alignat*}{4}
       &      &      x_2 & {}-{} &  4x_3 &= 8 \\
  2x_1 & {}-{} &  3x_2 & {}+{} &  2x_3 &= 1 \\
  4x_1 & {}-{} &  8x_2 & {}+{} & 12x_3 &= 1
\end{alignat*}

\begin{table}[H]
\centering
\renewcommand{\arraystretch}{1.2}

\begin{tabular}{|>{\centering\arraybackslash}m{4.5cm}|>{\centering\arraybackslash}m{6.5cm}|}
    \hline
    \textbf{Row operations} & \textbf{Augmented matrix} \\ \hline
    
    % --- Step 0 (Initial Matrix) ---
    \rule{0pt}{2.5em} 
    \textit{Initial Matrix}
    & 
    $\left[
    \begin{array}{ccc|c}
        0 &  1 & -4 & 8 \\
        2 & -3 &  2 & 1 \\
        4 & -8 & 12 & 1
    \end{array}
    \right]$ 
    \\[2em] \hline
    
    % --- Step 1 (Swap) ---
    \rule{0pt}{2.5em} 
    $R_1 \leftrightarrow R_2$ 
    & 
    $\left[
    \begin{array}{ccc|c}
        2 & -3 &  2 & 1 \\
        0 &  1 & -4 & 8 \\
        4 & -8 & 12 & 1
    \end{array}
    \right]$ 
    \\[2em] \hline
    
    % --- Step 2 ---
    \rule{0pt}{2.5em} 
    $R_3 - 2R_1 \rightarrow R_3$ 
    & 
    $\left[
    \begin{array}{ccc|c}
        2 & -3 &  2 &  1 \\
        0 &  1 & -4 &  8 \\
        0 & -2 &  8 & -1
    \end{array}
    \right]$ 
    \\[2em] \hline
    
    % --- Step 3 (Final Result) ---
    \rule{0pt}{2.5em} 
    $R_3 + 2R_2 \rightarrow R_3$ 
    & 
    $\left[
    \begin{array}{ccc|c}
        2 & -3 &  2 &  1 \\
        0 &  1 & -4 &  8 \\
        0 &  0 &  0 & 15
    \end{array}
    \right]$ 
    \\[2em] \hline
\end{tabular}
\end{table}
\noindent The last equation is
\[0x_1 + 0x_2 + 0x_3 = 15\]
i.e. $0=15$. Therefore the system is inconsistent.

\subsubsection{Example (Consistent system, infinite solution)}
Determine the solution to this linear system:
\begin{alignat*}{7}
    x_1  & {}+{} & 3x_2 & {}-{} & 2x_3 &       &       & {}+{} & 2x_5 &       &        &= 0 \\
   2x_1  & {}+{} & 6x_2 & {}-{} & 5x_3 & {}-{} &  2x_4 & {}+{} & 4x_5 & {}-{} &  3x_6  &= -1 \\
         &       &      &       & 5x_3 & {}+{} & 10x_4 &       &      & {}+{} & 15x_6  &= 5 \\
   2x_1  & {}+{} & 6x_2 &       &      & {}+{} &  8x_4 & {}+{} & 4x_5 & {}+{} & 18x_6  &= 6
\end{alignat*}
\begin{table}[H]
\centering
\renewcommand{\arraystretch}{1.2}

\begin{tabular}{|>{\centering\arraybackslash}m{3.5cm}|>{\centering\arraybackslash}m{10cm}|}
    \hline
    \textbf{Row operations} & \textbf{Augmented matrix} \\ \hline
    
    % --- Step 1: Initial Matrix ---
    \rule{0pt}{2.5em} 
    \textit{Initial System}
    & 
    $\left[
    \begin{array}{cccccc|c}
        1 & 3 & -2 & 0 & 2 & 0 & 0 \\
        2 & 6 & -5 & -2 & 4 & -3 & -1 \\
        0 & 0 & 5 & 10 & 0 & 15 & 5 \\
        2 & 6 & 0 & 8 & 4 & 18 & 6
    \end{array}
    \right]$ 
    \\[3em] \hline
    
    % --- Step 2 ---
    \rule{0pt}{2.5em} 
    $\begin{aligned}
        R_2 - 2R_1 &\rightarrow R_2 \\
        R_4 - 2R_1 &\rightarrow R_4
    \end{aligned}$ 
    & 
    $\left[
    \begin{array}{cccccc|c}
        1 & 3 & -2 & 0 & 2 & 0 & 0 \\
        0 & 0 & -1 & -2 & 0 & -3 & -1 \\
        0 & 0 & 5 & 10 & 0 & 15 & 5 \\
        0 & 0 & 4 & 8 & 0 & 18 & 6
    \end{array}
    \right]$ 
    \\[3em] \hline
    
    % --- Step 3 ---
    \rule{0pt}{2.5em} 
    $-R_2 \rightarrow R_2$ 
    & 
    $\left[
    \begin{array}{cccccc|c}
        1 & 3 & -2 & 0 & 2 & 0 & 0 \\
        0 & 0 & 1 & 2 & 0 & 3 & 1 \\
        0 & 0 & 5 & 10 & 0 & 15 & 5 \\
        0 & 0 & 4 & 8 & 0 & 18 & 6
    \end{array}
    \right]$ 
    \\[3em] \hline

    % --- Step 4 ---
    \rule{0pt}{2.5em} 
    $\begin{aligned}
        R_3 - 5R_2 &\rightarrow R_3 \\
        R_4 - 4R_2 &\rightarrow R_4
    \end{aligned}$ 
    & 
    $\left[
    \begin{array}{cccccc|c}
        1 & 3 & -2 & 0 & 2 & 0 & 0 \\
        0 & 0 & 1 & 2 & 0 & 3 & 1 \\
        0 & 0 & 0 & 0 & 0 & 0 & 0 \\
        0 & 0 & 0 & 0 & 0 & 6 & 2
    \end{array}
    \right]$ 
    \\[3em] \hline

    % --- Step 5 (Row Echelon Form Reached) ---
    \rule{0pt}{2.5em} 
    $\begin{aligned}
        R_3 &\leftrightarrow R_4 \\
        \tfrac{1}{6}R_3 &\rightarrow R_3
    \end{aligned}$ 
    & 
    $\left[
    \begin{array}{cccccc|c}
        1 & 3 & -2 & 0 & 2 & 0 & 0 \\
        0 & 0 & 1 & 2 & 0 & 3 & 1 \\
        0 & 0 & 0 & 0 & 0 & 1 & \frac{1}{3} \\
        0 & 0 & 0 & 0 & 0 & 0 & 0
    \end{array}
    \right]$ 
    \\[2.5em] % Add space, but NO \hline here
    
    % --- Note Row for REF ---
    & \textit{Note: Matrix is now in Row Echelon Form (REF)} \\[1em] \hline
    
     % --- Step 6 ---
    \rule{0pt}{2.5em} 
    $R_2 - 3R_3 \rightarrow R_2$ 
    & 
    $\left[
    \begin{array}{cccccc|c}
        1 & 3 & -2 & 0 & 2 & 0 & 0 \\
        0 & 0 & 1 & 2 & 0 & 0 & 0 \\
        0 & 0 & 0 & 0 & 0 & 1 & \frac{1}{3} \\
        0 & 0 & 0 & 0 & 0 & 0 & 0
    \end{array}
    \right]$ 
    \\[3em] \hline

    % --- Step 7 (Reduced Row Echelon Form Reached) ---
    \rule{0pt}{2.5em} 
    $R_1 + 2R_2 \rightarrow R_1$ 
    & 
    $\left[
    \begin{array}{cccccc|c}
        1 & 3 & 0 & 4 & 2 & 0 & 0 \\
        0 & 0 & 1 & 2 & 0 & 0 & 0 \\
        0 & 0 & 0 & 0 & 0 & 1 & \frac{1}{3} \\
        0 & 0 & 0 & 0 & 0 & 0 & 0
    \end{array}
    \right]$ 
    \\[2.5em] % Add space, but NO \hline here
    
    % --- Note Row for RREF ---
    & \textit{Note: Matrix is now in Reduced Row Echelon Form (RREF)} \\[1em] \hline
\end{tabular}
\end{table}
\noindent We first note that the system is consistent.

\noindent Also, not every variable column has a pivot (in either REF or RREF).
Hence, the system has infinitely many solutions.

\noindent From the RREF, we obtain this system of equations:
\begin{align*}
    x_1 + 3x_2 + 4x_4 + 2x_5 &= 0 \\
    x_3 + 2x_4 &= 0 \\
    x_6 &= \frac{1}{3}
\end{align*}

\noindent Solving for the leading variables in terms of the free variables,
\begin{align*}
    x_1 &= -3x_2 - 4x_4 - 2x_5 \\
    x_3 &= -2x_4 \\
    x_6 &= \frac{1}{3}
\end{align*}

\noindent The general solution is expressed parametrically,
\begin{align*}
    x_1 &= -3r - 4s - 2t \\
    x_2 &= r \\
    x_3 &= -2s \\
    x_4 &= s \\
    x_5 &= t \\
    x_6 &= \frac{1}{3}
\end{align*}

\noindent\textbf{Alternative}: Gaussian elimination with back-substitution for larger linear systems
(instead of Gauss-Jordan elimination).

\noindent From the REF,
\begin{align*}
    x_1 + 3x_2 - 2x_3 + 2x_5 &= 0 \\
    x_3 + 2x_4 + 3x_6 &= 0 \\
    x_6 &= \frac{1}{3}
\end{align*}

\noindent Solving for the leading variables,
\begin{align*}
    x_1 &= -3x_2 + 2x_3 - 2x_5 \\
    x_3 &= 1 - 2x_4 - 3x_6 \\
    x_6 &= \frac{1}{3}
\end{align*}

\noindent Back-substitute the bottom equations into those above it:
\begin{align*}
    x_1 &= -3x_2 - 4x_4 - 2x_5 \\
    x_3 &= -2x_4 \\
    x_6 &= \frac{1}{3}
\end{align*}

\noindent Assigning arbitary parameters $r,s,t$ to the free variables, we would get
the same general solution.

\subsubsection{Example (Consistent system, unique solution)}
Determine the solution to this linear system:
\begin{align*}
    2x_1 + x_2 - x_3 &= 8 \\
    -3x_1 - x_2 + 2x_3 &= -11 \\
    -2x_1 + x_2 + 2x_3 &= -3
\end{align*}
\noindent Row echelon form:
\[
\left[
\begin{array}{ccc|c}
    2 & 1 & -1 & 8 \\
    0 & \frac{1}{2} & \frac{1}{2} & 1 \\
    0 & 0 & -1 & 1
\end{array}
\right]
\]
\noindent First, note that the system is consistent.

\noindent Since every variable column contains a pivot,
there is a unique solution. (There are no free variables)

\medskip

\noindent Reduced row echelon form:
\[
\left[
\begin{array}{ccc|c}
    1 & 0 & 0 & 2 \\
    0 & 1 & 0 & 3 \\
    0 & 0 & 1 & -1
\end{array}
\right]
\]
\noindent Note that the RREF of every linear system with a unique solution
is the identity matrix.

\noindent The solution to the linear system can be read off directly from the RREF:
\[x_1=1,\:x_2=3,\:x_3=-1\]

\subsection{Linear combination of vectors}
\noindent Given vectors $\mathbf{v}_1, \mathbf{v}_2, \cdots, \mathbf{v}_p$ in $\mathbb{R}^n$ and given scalars $c_1, c_2, \dots, c_p$, the vector $\mathbf{y}$ defined by
\[
\mathbf{y} = c_1\mathbf{v}_1 + \dots + c_p\mathbf{v}_p
\]
is called a linear combination of $\mathbf{v}_1, \mathbf{v}_2, \dots, \mathbf{v}_p$ with \textbf{weights} $c_1, c_2, \dots, c_p$.

\subsubsection{Example} 
Let $\mathbf{a}_1 = \begin{bmatrix} 1 \\ -2 \\ -5 \end{bmatrix}$, 
$\mathbf{a}_2 = \begin{bmatrix} 2 \\ 5 \\ 6 \end{bmatrix}$, and 
$\mathbf{b} = \begin{bmatrix} 7 \\ 4 \\ -3 \end{bmatrix}$. 
Determine whether $\mathbf{b}$ can be written as a linear combination of $\mathbf{a}_1$ and $\mathbf{a}_2$.

\noindent The vector equation $x_1\mathbf{a}_1 + x_2\mathbf{a}_2 = \mathbf{b}$ is represented by the augmented matrix $[ \mathbf{a}_1 \ \mathbf{a}_2 \mid \mathbf{b} ]$:
\[
\left[
\begin{array}{cc|c}
    1 & 2 & 7 \\
    -2 & 5 & 4 \\
    -5 & 6 & -3
\end{array}
\right]
\]

After row reduction, we obtain the Reduced Row Echelon Form (RREF):
\[
\left[
\begin{array}{cc|c}
    1 & 0 & 3 \\
    0 & 1 & 2 \\
    0 & 0 & 0
\end{array}
\right]
\]

From the RREF, we can read the unique solution for the weights:
\[
x_1 = 3, \quad x_2 = 2
\]

Thus, $\mathbf{b}$ is a linear combination of $\mathbf{a}_1$ and $\mathbf{a}_2$:
\[
\mathbf{b} = 3\mathbf{a}_1 + 2\mathbf{a}_2
\]

\subsection{Span}
If $\mathbf{v}_1,\mathbf{v}_2,\ldots,\mathbf{v}_p$ are vectors in $\mathbb{R}^n$,
then Span$\{\mathbf{v}_1,\mathbf{v}_2,\ldots,\mathbf{v}_p\}$ is the set of all vectors that can be
written in the form
\[c_1\mathbf{v}_1 + c_2\mathbf{v}_2 + \ldots + c_p\mathbf{v}_p\]
where $c_1,c_2,\ldots,c_p$ are scalars.

\noindent It is also called the subset of $\mathbb{R}^n$ spanned by $\mathbf{v}_1,\mathbf{v}_2,\ldots,\mathbf{v}_n$.

\subsubsection{Geometric intepretation}
\begin{itemize}
    \item Span of a single nonzero vector:
    All multiples of that vector form a line passing through the origin.
    \item Span of two non-parallel vectors:
    All combinations of these two vectors form a plane passing through the origin.
    \item Span of three vectors:
    If they are linearly independent to each other, their span is $\mathbb{R}^3$.
\end{itemize}

\begin{figure}[htbp]
    \centering
    \includegraphics[width=0.5\textwidth]{images/span.png} % Replace 'filename' with your image name
    \caption*{Spans as lines, planes, and real space}
    \label{fig:span geometric interpretation}
\end{figure}

\subsubsection{Example}
Let $\mathbf{a}_1 = \begin{bmatrix} 1 \\ -2 \\ 3 \end{bmatrix}$, $\mathbf{a}_2 = \begin{bmatrix} 5 \\ -13 \\ -3 \end{bmatrix}$, and $\mathbf{b} = \begin{bmatrix} -3 \\ 8 \\ 1 \end{bmatrix}$. Determine if $\mathbf{b}$ is in the $\text{Span}\{\mathbf{a}_1, \mathbf{a}_2\}$.

\medskip

\noindent The augmented matrix is:
\[
\left[
\begin{array}{cc|c}
1 & 5 & -3 \\
-2 & -13 & 8 \\
3 & -3 & 1
\end{array}
\right]
\]

\noindent Performing row reduction to obtain row echelon form,
\[
\left[
\begin{array}{cc|c}
1 & 5 & -3 \\
0 & -3 & 2 \\
0 & 0 & -2
\end{array}
\right]
\]

\noindent The third row corresponds to the equation $0x_1 + 0x_2 = -2$, which is a contradiction.

\noindent Therefore, the system is inconsistent, and $\mathbf{b}$ is not in the $\text{Span}\{\mathbf{a}_1, \mathbf{a}_2\}$.

\subsection{The Matrix Equation $A\mathbf{x} = \mathbf{b}$}
Recall matrix multiplication as linear combinations of columns of $A$.
\begin{figure}[H]
    \centering
    \includegraphics[width=0.8\textwidth]{images/linearcomb.png}
    % \caption*{Linear combination of column vectors of $A$, where the coefficients
    % are elements of the single column vector.}
    \label{fig:matrix_multiplication}
\end{figure}

\noindent If $A$ is an $m\times n$ matrix, with columns $a_1, a_2,\ldots,a_n$,
and if $\mathbf{x}$ is in $\mathbb{R}^n$, then $A\mathbf{x}$ is the linear combination
of the columns of A using the corresponding entries in $\mathbf{x}$ as weights,
\[A\mathbf{x} = \begin{bmatrix}
    \mathbf{a}_1 & \mathbf{a}_1 & \cdots & \mathbf{a}_n \\
\end{bmatrix} \begin{bmatrix}
    x_1 \\
    \vdots \\
    x_n
\end{bmatrix} = x_1\mathbf{a}_1 + x_1\mathbf{a}_1 + \ldots + x_n\mathbf{a}_n\]

\noindent Further, if $\mathbf{b}$ is in $\mathbb{R}^m$, then
\[A\mathbf{x} = \mathbf{b}\]
has the same solution as 
\[x_1\mathbf{a}_1 + x_1\mathbf{a}_1 + \ldots + x_n\mathbf{a}_n = \mathbf{b}\]
i.e. the same solution to the linear system represented by this augmented matrix
\[\begin{bmatrix}
    \mathbf{a}_1 & \mathbf{a}_2 & \cdots & \mathbf{a}_n & \mathbf{b}
\end{bmatrix}\]

\subsubsection{Example}
\noindent Let $A = 
\begin{bmatrix} 
1 & 3 & 4 \\ 
-4 & 2 & -6 \\ 
-3 & -2 & -7 
\end{bmatrix}$ 
and $\mathbf{b} = 
\begin{bmatrix} 
b_1 \\ 
b_2 \\ 
b_3 
\end{bmatrix}$. 
Is the equation $A\mathbf{x} = \mathbf{b}$ consistent for all possible values of $b_1, b_2, b_3$?

\noindent The augmented matrix is
\[ 
\left[
\begin{array}{ccc|c}
1 & 3 & 4 & b_1 \\
-4 & 2 & -6 & b_2 \\
-3 & -2 & -7 & b_3
\end{array}
\right]
\]

\noindent Performing row reduction, the row echelon form is
\[
\left[
\begin{array}{ccc|c}
1 & 3 & 4 & b_1 \\
0 & 14 & 10 & 4b_1 + b_2 \\
0 & 0 & 0 & b_1 - \frac{1}{2}b_2 + b_3
\end{array}
\right]
\]
For the system to be consistent, we need
\[
b_1 - \frac{1}{2}b_2 + b_3 = 0
\]
The columns of A lie on a plane passing through the origin in $\mathbb{R}^3$,
and the equation of the plane is $2x-y+2x=0$.

\subsubsection{Spanning $\mathbb{R}^m$}
If Span\{$\mathbf{v}_1,\mathbf{v}_2,\ldots,\mathbf{v}_p$\}=$\mathbb{R}^m$, every vector in $\mathbb{R}^m$
is a linear combination of $\mathbf{v}_1,\ldots,\mathbf{v}_p$.

\medskip

\noindent These statements are logically equivalent:
\begin{enumerate}
    \item $A\mathbf{x} = \mathbf{b}$ has a solution for each $\mathbf{b}$ in $\mathbb{R}^m$.\newline
    i.e. Any vector in $\mathbb{R}^m$ can be found through $A$.
    \item Each $\mathbf{b}$ in $\mathbf{R}^m$ is a linear combination of the columns of $A$.
    \item The columns of $A$ span $\mathbb{R}^m$.
    \item $A$ has a pivot position in every row.\newline
    If not, the last row will have all zero entries. Then $\mathbf{b}$ can be constructed with a 1 in the last row,
    making the linear system inconsistent.\newline
    \textbf{**}Not to be confused with a pivot in every variable column, which shows uniqueness of solution for the linear system.
\end{enumerate}

\noindent Some properties, for $\mathbf{u}$ and $\mathbf{v}$ in $\mathbb{R}^m$, and $c$ a scalar:
\begin{enumerate}
    \item $A(\mathbf{u} + \mathbf{v}) = A\mathbf{u} + A\mathbf{v}$
    \item $A(c\mathbf{u}) = c(A\mathbf{u})$
\end{enumerate}

\subsubsection{Note on Pivots}
A pivot in every column $\iff$ there are no free variables $\iff$ the linear system $A\mathbf{x}=\mathbf{b}$
has at most one solution. $\iff$ Columns of $A$ are linearly independent. The system is either inconsistent or has a unique solution.
% \begin{equation*}
% \left[
% \begin{array}{cc|c}
% 1 & 0 & 1 \\
% 0 & 1 & 1 \\
% 0 & 0 & 1
% \end{array}
% \right]
% \end{equation*}
% \begin{equation*}
%     \left[
%     \begin{array}{cc|c}
%     1 & 2 & 5 \\
%     0 & 1 & 2 \\
%     0 & 0 & 0
%     \end{array}
%     \right]
% \end{equation*}

\begin{figure}[h!]
     \centering
     % First Matrix
     \begin{subfigure}[b]{0.45\textwidth}
         \centering
         $\left[
         \begin{array}{cc|c}
         \boxed{1} & 2 & 5 \\
         0 & \boxed{1} & 2 \\
         0 & 0 & 0
         \end{array}
         \right]$
         \caption{Unique Solution.}
         \label{fig:unique}
     \end{subfigure}
     \hfill
     % Second Matrix
     \begin{subfigure}[b]{0.45\textwidth}
         \centering
         $\left[
         \begin{array}{cc|c}
         \boxed{1} & 2 & 5 \\
         0 & \boxed{1} & 2 \\
         0 & 0 & 1
         \end{array}
         \right]$
         \caption{No Solution ($0=1$)}
         \label{fig:no_sol}
     \end{subfigure}
     
     \caption*{Pivots in every column.}
     \label{fig:matrix_comparison}
\end{figure}

\noindent A pivot in every row $\iff$ the linear system $A\mathbf{x}=\mathbf{b}$
has at least one solution, for every $\mathbf{b}$. This also means the columns of $A$
span $\mathbb{R}^m$, where $m$ is the number of rows.

\noindent $m$ pivots mean $m$ dimensions ($m$ linearly independent rows) to span $\mathbb{R}^m$.

\begin{figure}[H]
     \centering
     % Wide Matrix
     \begin{subfigure}[b]{0.48\textwidth}
         \centering
         $\left[
         \begin{array}{ccc|c}
         \boxed{1} & 0 & 2 & b_1 \\
         0 & \boxed{1} & 3 & b_2
         \end{array}
         \right]$
         \caption{Wide matrices with pivot in every row are always consistent and have
          infinitely many solutions (at least one free variable left over).}
         \label{fig:wide_matrix}
     \end{subfigure}
     \hfill
     % Tall Matrix
     \begin{subfigure}[b]{0.48\textwidth}
         \centering
         $\left[
         \begin{array}{cc|c}
         \boxed{1} & 0 & b_1 \\
         0 & \boxed{1} & b_2 \\
         0 & 0 & b_3
         \end{array}
         \right]$
         \caption{Tall matrices cannot have a pivot in every row. Row of zeros leads to "No Solution" if $b_3 \neq 0$.}
         \label{fig:tall_matrix}
     \end{subfigure}
     
     \caption*{Pivots on every row.}
\end{figure}

\noindent There is a pivot in every row and column $\iff$ $A$ is a square matrix $\iff$ 
The system $A\mathbf{x}=\mathbf{b}$ has a unique solution for any $\mathbf{b}$.

\begin{figure}[h!]
    \centering
    $\left[
    \begin{array}{ccc|c}
    \boxed{1} & 0 & 0 & 5 \\
    0 & \boxed{1} & 0 & -2 \\
    0 & 0 & \boxed{1} & 4
    \end{array}
    \right]$
    \caption*{Pivots in every row and column lead to a unique solution.}
    \label{fig:square_unique}
\end{figure}

\medskip
\noindent Hence, wide matrix can never have unique solution,
tall matrix can never have solution for every b.

\subsection{Solution Sets of Linear Systems}
Parametric vector form of the solution set of $A\mathbf{x}=\mathbf{b}$ is
\[\mathbf{x} = \mathbf{p} + t\mathbf{v},\quad t\in\mathbb{R}\]

\noindent If $\mathbf{b}=\mathbf{0}$, it is a homogenous system of equations.

\noindent Parametric vector form for the solution to a homogenous system is
\[\mathbf{x} = t\mathbf{v}\]

\noindent $\mathbf{p}$ is then one particular solution to $A\mathbf{x}=\mathbf{b}$
(corresponding to t=0).

\noindent $\mathbf{p}$ basically translates $t\mathbf{v}$ to pass through $\mathbf{p}$ instead of the origin.    

\subsubsection{Example (homogenous system)}
Determine if this homogenous system has a nontrivial solution:
\[
\begin{aligned}
3x_1 + 5x_2 - 4x_3 &= 0 \\
-3x_1 - 2x_2 + 4x_3 &= 0 \\
6x_1 + x_2 - 8x_3 &= 0
\end{aligned}
\]
The augmented matrix is
\[
\left[
\begin{array}{ccc|c}
3 & 5 & -4 & 0 \\
-3 & -2 & 4 & 0 \\
6 & 1 & -8 & 0
\end{array}
\right]
\]
Performing row reduction operations, the row echelon form is
\[
\left[
\begin{array}{ccc|c}
3 & 5 & -4 & 0 \\
0 & 3 & 0 & 0 \\
0 & 0 & 0 & 0
\end{array}
\right]
\]
We get
\[
\begin{aligned}
3x_1 + 5x_2 - 4x_3 &= 0 \\
3x_2 &= 0 
\end{aligned}
\]
Solving,
\[
3x_1 = 4x_3 \implies x_1 = \frac{4}{3}x_3
\]
Since $x_3$ is a free variable, we can express the solution vector $\mathbf{x}$ as:
\[
\mathbf{x} = 
\begin{bmatrix} 
x_1 \\ 
x_2 \\ 
x_3 
\end{bmatrix} 
= 
\begin{bmatrix} 
\frac{4}{3}x_3 \\ 
0 \\ 
x_3 
\end{bmatrix}
= x_3 
\begin{bmatrix} 
\frac{4}{3} \\ 
0 \\ 
1 
\end{bmatrix}
\]
Every solution to $A\mathbf{x}=\mathbf{0}$ in this case is a scalar multiple 
of $\mathbf{v}=\begin{bmatrix}
    \frac{4}{3} \\ 0 \\ 1
\end{bmatrix}$

\subsubsection{Example (non-homogenous system)}
Describe all solutions of this non-homogenous system:
\[
\begin{aligned}
3x_1 + 5x_2 - 4x_3 &= 7 \\
-3x_1 - 2x_2 + 4x_3 &= -1 \\
6x_1 + x_2 - 8x_3 &= 4
\end{aligned}
\]
The augmented matrix is
\[
\left[
\begin{array}{ccc|c}
3 & 5 & -4 & 7 \\
-3 & -2 & 4 & -1 \\
6 & 1 & -8 & -4
\end{array}
\right]
\]
The reduced row echelon form is
\[
\left[
\begin{array}{ccc|c}
1 & 0 & -\frac{4}{3} & -1 \\
0 & 1 & 0 & 2 \\
0 & 0 & 0 & 0
\end{array}
\right]
\]
We get
\[
\left\{
\begin{aligned}
x_1 - \frac{4}{3}x_3 &= -1 \\
x_2 &= 2 \\
x_3 &\text{ is free}
\end{aligned}
\right.
\]
The solution set is
\[
\mathbf{x} = 
\begin{bmatrix} 
x_1 \\ 
x_2 \\ 
x_3 
\end{bmatrix} 
= 
\begin{bmatrix} 
-1 \\ 
2 \\ 
0 
\end{bmatrix} 
+ t 
\begin{bmatrix} 
\frac{4}{3} \\ 
0 \\ 
1 
\end{bmatrix},\quad
t\in\mathbb{R}
\]

\subsection{Linear Independence}
A indexed set of vectors $\{\mathbf{v}_1,\ldots,\mathbf{v}_p\}$ in $\mathbb{R}^n$
is linearly independent iff the vector equation
\[x_1\mathbf{v}_1 + x_2\mathbf{v}_2 + \cdots + x_p\mathbf{v}_p = \mathbf{0}\]
has only the trivial solution.

\noindent If $\exists \text{ weights } c_1,\ldots,c_p$ that are not all zero, s.t.
\[c_1\mathbf{v}_1 + c_2\mathbf{v}_2 +\cdots+ c_p\mathbf{v}_p = \mathbf{0}\]
then the set of vectors is linearly dependent.

\noindent In other words, if $S = \{\mathbf{v}_1,\cdots,\mathbf{v}_p\}$ is linearly dependent,
and $\mathbf{v}_1 \neq 0$, then some $\mathbf{v}_j$, $j>1$, is a linear combination of 
the preceding vectors $\mathbf{v}_1,\cdots,\mathbf{v}_{j-1}$.

\smallskip
\noindent In a matrix, if all the columns have a pivot, then the columns are linearly independent.

\subsubsection{Example}

Consider the set of vectors in $\mathbb{R}^3$:
\[
\mathbf{v}_1 = \begin{bmatrix} 1 \\ 1 \\ 0 \end{bmatrix}, \quad 
\mathbf{v}_2 = \begin{bmatrix} 1 \\ 0 \\ 1 \end{bmatrix}, \quad 
\mathbf{v}_3 = \begin{bmatrix} 3 \\ 1 \\ 2 \end{bmatrix}, \quad 
\mathbf{v}_4 = \begin{bmatrix} 0 \\ 1 \\ 1 \end{bmatrix}
\]

\noindent First construct the matrix $A = [\mathbf{v}_1 \ \mathbf{v}_2 \ \mathbf{v}_3 \ \mathbf{v}_4]$ and reduce it to REF to locate the pivots:
\[
A \sim \begin{bmatrix} 
\mathbf{1} & 1 & 3 & 0 \\ 
0 & \mathbf{1} & 2 & -1 \\ 
0 & 0 & 0 & \mathbf{1} 
\end{bmatrix}
\]
Column 1 and Column 2 have pivots, Column 3 is the first column without a pivot.
This indicates that $\mathbf{v}_3$ is the first vector in the sequence that is linearly dependent on those preceding it.

\noindent We can extract the vectors up to $\mathbf{v}_3$ into an augmented matrix $[\mathbf{v}_1 \ \mathbf{v}_2 \mid \mathbf{v}_3]$:
\[
\left[ \begin{array}{cc|c} 
\mathbf{1} & 1 & 3 \\ 
0 & \mathbf{1} & 2 \\ 
0 & 0 & 0 
\end{array} \right]
\]
\begin{itemize}
    \item Consistency: The absence of a pivot in the augmented column (the right-most column) proves the system is consistent; thus, a linear combination exists.
    \item Unique Weights: From the REF, we can back-substitute:
    \begin{align*}
        x_2 &= 2 \\
        x_1 + x_2 = 3 \implies x_1 + 2 &= 3 \implies x_1 = 1
    \end{align*}
\end{itemize}

\noindent Since $\mathbf{v}_3$ corresponds to a non-pivot column in the original set, it is redundant. Specifically:
\[ \mathbf{v}_3 = 1\mathbf{v}_1 + 2\mathbf{v}_2 \]

\noindent Takeaway: If there are columns with no pivot, the set of vectors is linearly dependent.

\subsubsection{Wide matrices}
Naturally, columns of a wide matrix are linearly dependent.
\[
n
\begin{array}{c}
  \scriptstyle p \\
  \left[ 
    \begin{array}{ccccc}
      * & * & * & * & * \\
      * & * & * & * & * \\
      * & * & * & * & * \end{array} 
  \right]
\end{array}
\]
If $p>n$, the matrix cannot have a pivot in every column $\Rightarrow$
there must be a free variable $\Rightarrow$ $A\mathbf{x} = \mathbf{0}$
has non trivial solution $\Rightarrow$ columns of $A$ are linerly dependent.

\subsubsection{Set with zero vector}
If $S  = \{\mathbf{v}_1,\ldots,\mathbf{v}_p\}$ in $\mathbb{R}^n$ contains
the zero vector, then $S$ is linearly dependent.

\smallskip
\noindent WLOG, suppose $\mathbf{v}_1 = \mathbf{0}$, then
\[1\mathbf{v}_1 + 0\mathbf{v}_2 +\cdots +0\mathbf{v}_p = \mathbf{0}\]
shows that $S$ is linearly dependent.

\subsection{Linear Transformations}
\noindent A transformation $T : \mathbb{R}^n \to \mathbb{R}^m$ is a rule that 
assigns to each vector $\mathbf{x}$ in $\mathbb{R}^n$ a vector $T(\mathbf{x})$ in $\mathbb{R}^m$.
\begin{itemize}
    \item $\mathbb{R}^n$ - domain of $T$
    \item $\mathbb{R}^m$ - codomain of $T$
    \item $T(\mathbf{x})$ - image of $\mathbf{x}$
    \item Set of all images $T(\mathbf{x})$ - range of $T$
\end{itemize}


\noindent A transformation $T$ is linear if it satisfies:
\begin{enumerate}
    \item $T(\mathbf{u} + \mathbf{v}) = T(\mathbf{u}) + T(\mathbf{v})$
    \item $T(c\mathbf{u}) = cT(\mathbf{u})$
\end{enumerate}
\noindent Generalization:
\[
T(c_1 \mathbf{v}_1 + \dots + c_p \mathbf{v}_p) = c_1 T(\mathbf{v}_1) + \dots + c_p T(\mathbf{v}_p)
\]

\noindent Example:
$T(\mathbf{x}) = r\mathbf{x}$, $T : \mathbb{R}^2 \to \mathbb{R}^2$ 
is called a \textbf{contraction} when $0 \leq r \leq 1$ and a \textbf{dilation} when $r > 1$.

\subsubsection{Finding the Standard Matrix}
Let $T:\mathbb{R}^n\rightarrow\mathbb{R}^m$ be a linear transformation.
Then there exists a unique matrix $A$ such that 
\[T(\mathbf{x}) = A\mathbf{x} \text{ for all } \mathbf{x} \text{ in } \mathbb{R}^n\]
$A$ is also known as the standard matrix for the linear transformation $T$.

\smallskip
\noindent Here, $A$ is the $m\times n$ matrix whose $j^{\text{th}}$ column is the 
vector $T(e_j)$, where $e_j$ is the $j^{\text{th}}$ column of the 
identity matrix in $\mathbb{R}^n$, i.e.
\[ A = \begin{bmatrix} T(\mathbf{e}_1) & T(\mathbf{e}_2) & \dots & T(\mathbf{e}_n) \end{bmatrix} \]

\noindent For $T:\mathbb{R}^n\rightarrow\mathbb{R}^m$, $T(\mathbf{x})=A\mathbf{x}$,
\[A \text{ is an } m\times n \text{ matrix.}\]

\bigskip
\noindent
\begin{minipage}{\textwidth}
\textbf{Example:}

\noindent Let $I_2 = \begin{bmatrix} 1 & 0 \\ 0 & 1 \end{bmatrix}$, 
so $\mathbf{e}_1 = \begin{bmatrix} 1 \\ 0 \end{bmatrix}$ and $\mathbf{e}_2 = \begin{bmatrix} 0 \\ 1 \end{bmatrix}$. 
Suppose $T$ is a linear transformation from $\mathbb{R}^2$ to $\mathbb{R}^3$ 
such that $T(\mathbf{e}_1) = \begin{bmatrix} 5 \\ -7 \\ 2 \end{bmatrix}$ and $T(\mathbf{e}_2) = \begin{bmatrix} -3 \\ 8 \\ 0 \end{bmatrix}$. 
Find the formula for the image of an arbitrary $\mathbf{x}$ in $\mathbb{R}^2$.

\vspace{1em}

\noindent There exists a unique matrix $A$ such that $T(\mathbf{x}) = A\mathbf{x}$ for all $\mathbf{x}$ in $\mathbb{R}^n$. 
The columns of $A$ are the images of the basis vectors $\mathbf{e}_1$ and $\mathbf{e}_2$:
\[
A = [T(\mathbf{e}_1) \quad T(\mathbf{e}_2)] = \begin{bmatrix} 5 & -3 \\ -7 & 8 \\ 2 & 0 \end{bmatrix}
\]
For an arbitrary $\mathbf{x} = \begin{bmatrix} x_1 \\ x_2 \end{bmatrix}$ in $\mathbb{R}^2$, the formula for the image is:
\[
T(\mathbf{x}) = \begin{bmatrix} 5 & -3 \\ -7 & 8 \\ 2 & 0 \end{bmatrix} \begin{bmatrix} x_1 \\ x_2 \end{bmatrix} = \begin{bmatrix} 5x_1 - 3x_2 \\ -7x_1 + 8x_2 \\ 2x_1 \end{bmatrix}
\]
\end{minipage}

\noindent \textbf{Example:}

\noindent Let $T : \mathbb{R}^2 \to \mathbb{R}^2$ be the transformation that rotates each point in $\mathbb{R}^2$ about the origin through an angle $\phi$, with counterclockwise rotation for a positive angle. Find the standard matrix $A$ for this transformation.

\vspace{1em}
\noindent To find the standard matrix $A$, we determine the images of the standard basis vectors $\mathbf{e}_1 = \begin{bmatrix} 1 \\ 0 \end{bmatrix}$ and $\mathbf{e}_2 = \begin{bmatrix} 0 \\ 1 \end{bmatrix}$ under the rotation $\phi$:

\begin{itemize}
    \item The vector $\mathbf{e}_1$ rotates to $T(\mathbf{e}_1) = \begin{bmatrix} \cos \phi \\ \sin \phi \end{bmatrix}$.
    \item The vector $\mathbf{e}_2$ rotates to $T(\mathbf{e}_2) = \begin{bmatrix} -\sin \phi \\ \cos \phi \end{bmatrix}$.
\end{itemize}

The standard matrix $A$ is formed by these columns:
\[
A = [T(\mathbf{e}_1) \quad T(\mathbf{e}_2)] = \begin{bmatrix} \cos \phi & -\sin \phi \\ \sin \phi & \cos \phi \end{bmatrix}
\]

\end{document}